% Historical drafting fragment. The integrated manuscript is docs/paper/review-2026-10-02/security.tex and supplement-security.tex.
% Integration draft, 2026-10-02. Based on Claude's second review, audited
% against 7a2517e plus the unused-input-certificate check in this review.
% This is a section fragment, not a standalone compiled paper.
% Requires lemma/proposition environments and IEEEproof; reconcile labels
% with the main manuscript before inclusion. No new benchmark claims.

\subsection{Source Recovery and Decision-First Repair}
\label{sec:c3-security-draft}

\paragraph{State and assumptions.}
For an output $x$, let $o_x$ denote its direct obligation and
$\mathcal D(x)$ its immutable compensation decision. In reachable states
from a valid initial configuration, a decision exists exactly when the
obligation is \textsf{Repaired} or \textsf{Recovered}. Here,
\textsf{Repaired} denotes economic compensation, not physical installation
of revised historical bytes. Let $\mathcal R$ contain representation
revisions, tasks, batch references, and installation metadata, and let
$\pi_E$ retain public balances, spend states, obligations, coverage, grant
usage, and fee accounting while excluding representation metadata.

We use the fixed-membership fault model, atomic transitions, signature
unforgeability, and fact binding defined in the system model. Honest members
retain the payment, admission vector, and direct input certificates recorded
before signing. Correct committee replicas retain identical original blocks
and original block metadata required by outstanding obligations,
representation, and replay. Public verification binds each input certificate
to its network, organization, fact, and claimed output, rejects duplicates,
and requires an exact correspondence with certificate inputs. The admission
vector contains one CAL entry bound to the issuer. These properties are
enforced by the protocol checks; the arguments below reason over their
transition semantics.

\begin{lemma}[Recovery preserves authorization]
\label{lem:c3-recovery-authority}
A recovery job created by an honest member resubmits an already authorized
payment. It requires no new owner or member signature and does not reserve
additional signing capacity. Once a direct obligation is public, at least
two honest original signers retain the corresponding approved material.
\end{lemma}
\begin{IEEEproof}
The member indexes its stored approval by the fact in the public input QC.
The stored payment and admission vector must produce that same fact. The
member combines this material and its retained direct input certificates
with the exposed QC, then applies ordinary payment verification before
creating an outbox entry. Fact binding and quorum verification prevent a
different payment core from satisfying these checks. The transition writes
only the outbox: it neither signs nor changes the original approval, input
locks, or grant reservation. Existing outbox and observed-execution marks
suppress duplicate jobs.

An obligation can arise only from a successfully executed payment consuming
a certified missing input. Its QC is therefore public. Three distinct
signers include at least two honest members, each of which stored the
approved material before returning its signature. Thus withholding the
complete certificate at the original client alone does not remove the
materials needed for recovery. Actual execution additionally requires a
continuing honest follower, fair delivery, and the source payment to satisfy
the ordinary execution rules.
\end{IEEEproof}

\begin{lemma}[Compensation takes effect at most once]
\label{lem:c3-decision-once}
Each direct obligation incurs at most one reserve debit and its associated
economic effects. Repeating the same decision has no effect, including
after repayment; a conflicting decision is rejected.
\end{lemma}
\begin{IEEEproof}
A decision first checks the immutable record $\mathcal D(x)$. An exact
match returns an unapplied result with no state changes, economic effects,
or fee outputs; a mismatch is rejected. In the absence of a decision, the
transition requires an \textsf{Open} obligation, a reached public deadline,
and a matching historical position, amount, and funding commitment. It
atomically debits the issuer reserve, closes the gap, updates fees and
coverage, and records $\mathcal D(x)$.

If the source executes first, the obligation is no longer \textsf{Open} and
cannot be compensated. If compensation occurs first, subsequent source
execution returns the corresponding funds to the issuer through the single
\textsf{Repaired}$\rightarrow$\textsf{Recovered} transition. Repayment does
not erase $\mathcal D(x)$ or create another user output. Induction over the
public command order therefore gives at-most-once compensation and
repayment.
\end{IEEEproof}

\begin{lemma}[Representation has no economic effect]
\label{lem:c3-representation-inert}
For every successful single or batch representation transition $r$,
$\pi_E(r(S))=\pi_E(S)$. Legal physical installation likewise leaves this
economic projection unchanged. An effective representation requires a
committed compensation decision for every modified input slot.
\end{lemma}
\begin{IEEEproof}
Representation checks the immutable decision, current canonical target,
base revision, preceding digest, and permitted replacement. Its write set
contains only revisions, representation commands, pending-index removals,
and task metadata. It does not invoke the accounting transition. Followers
map the representation result to an empty economic effect. Installation
verifies the committed task and version before updating historical storage;
it does not execute payments or compensation again. The write sets prove
the projection equality for one step and, by induction, for any sequence
of such steps.
\end{IEEEproof}

The CAL conservation argument therefore uses the public decision as its
compensation event: for amount $a$, $\Delta A=\Delta D_{\mathrm{open}}=-a$
and $\Delta U=0$, preserving $U+A-D_{\mathrm{open}}$. Source resubmission
does not release signing capacity; representation and installation have
zero economic increment. The existing coverage argument applies to the
same economic transitions. Cross-history comparisons fix the economic
commands, public times, and historical target coordinates; projection
equality does not identify entire application states or permit arbitrary
deletion and renumbering of blocks.

\begin{lemma}[Identity preservation and independent installation]
\label{lem:c3-identity}
An effective representation preserves existing owner authorizations,
transaction and output identities, successor references, and the complete
original BlockID. Installing it does not require successor payments to be
re-signed or re-executed.
\end{lemma}
\begin{IEEEproof}
The executor constructs the canonical body by changing only the declared
funding slots and their openings. Fixed authorization fields and input
commitments remain unchanged, and replacements retain the encoded width.
It verifies the new openings, recomputes the revised body's block hash and
part-set header, and accepts only equality with the original BlockID.
Base and previous-digest checks serialize competing canonical revisions.
The committed task binds the permitted bytes. An installer advances an
older local version, checks an equal version, and treats a newer committed
version as already covering the task. Consequently physical progress
cannot revise an economic decision or invalidate an unchanged successor
reference.
\end{IEEEproof}

For a fixed original body and fixed set of decided slots, representation
also has a conditional convergence property: if the chameleon relation has
a unique valid opening for each fixed message and commitment, the final
body and part openings are independent of grouping and order. Each slot's
replacement is determined by its decision; part contexts do not depend on
the revision counter. Revision numbers, command histories, and application
hashes need not coincide across different groupings. This statement uses
opening uniqueness under valid cryptographic setup, separately from the
authorization and at-most-once arguments above.

\begin{proposition}[Economic closure does not wait for adaptation]
\label{prop:c3-closure}
An outstanding obligation can be closed by its valid public compensation
decision once its public deadline is reached and the decision is included,
provided the required original history, reserve, and fee escrow are
available. This transition does not require an adaptation contribution.
Representation completion separately requires valid adaptation material,
a current canonical revision, and progress of public submission and local
installation.
\end{proposition}
\begin{IEEEproof}
The decision path validates the obligation and historical slot and performs
accounting without invoking an adaptation endpoint. The admission rule
reserves a bounded compensation charge for each certified input, and the
coverage argument bounds reserve obligations under its stated assumptions.
Unavailable adaptation therefore prevents neither the economic decision
nor later repayment. In the configured threshold service, three valid
contributions suffice; three responding honest holders ensure this even
when the fourth withholds its contribution. Shared execution resources and
public ordering still govern latency, so this separates protocol
dependencies without asserting an unconditional completion-time bound.
\end{IEEEproof}
